\documentclass{article}
\usepackage{polski}
\usepackage[table]{xcolor}
\usepackage[a4paper, margin=1in]{geometry}
%usepackage{array} % Wymagane dla justowania tekstu w pionie w kolumnach 'm'

\title{Model biznesowy: Zarządzanie wyposażeniem}
\author{}
\date{}

\begin{document}
    \maketitle

    % --- PRZYPADEK 4: Zakup nowego wyposażenia ---
    \begin{tabular}{|c|c|}%{|m{10em}|m{20em}|}
        \hline
        \textbf{Nazwa} & Zakup nowego wyposażenia \\
        \hline
        \textbf{Numer PU} & 4 \\
        \hline
        \textbf{Aktorzy} & Administrator, Dostawca \\
        \hline
        \textbf{Krótki opis} & Proces polegający na wybraniu, zamówieniu i wprowadzeniu do systemu nowego sprzętu w celu uzupełnienia zasobów. \\
        \hline
        \textbf{Warunki wstępne} & Istnieje zapotrzebowanie na nowy sprzęt. \\
        \hline
        \textbf{Warunki końcowe} & Nowe wyposażenie znajduje się w bazie danych systemu ze statusem „Nowe”. \\
        \hline
        \textbf{Główny przepływ zdarzeń} & 
            1. Administrator wybiera opcję „Dodaj nowe wyposażenie”. \newline
            2. Administrator wprowadza dane sprzętu (nazwa, producent, model, cena). \newline
            3. Administrator zatwierdza wprowadzone dane. \newline
            4. System generuje unikalny numer inwentarzowy. \newline
            5. System zapisuje dane w bazie. \\
        \hline
        \textbf{Alternatywne przepływy zdarzeń} & 
            2a. Wprowadzone dane są niepoprawne – system wyświetla komunikat o błędzie. \newline
            3a. Anulowanie zakupu – Administrator przerywa procedurę, dane nie są zapisywane. \\
        \hline
        \textbf{Specjalne wymagania} & Interfejs musi pozwalać na szybkie wprowadzanie wielu sztuk sprzętu tego samego typu. \\
        \hline
        \textbf{Notatki} & brak \\
        \hline
    \end{tabular}

    \vspace{1cm} % Odstęp między tabelami

    % --- PRZYPADEK 5: Instalacja nowego wyposażenia ---
    \begin{tabular}{|c|c|}%{|m{10em}|m{20em}|}
        \hline
        \textbf{Nazwa} & Instalacja nowego wyposażenia \\
        \hline
        \textbf{Numer PU} & 5 \\
        \hline
        \textbf{Aktorzy} & Technik \\
        \hline
        \textbf{Krótki opis} & Fizyczne zamontowanie lub ustawienie sprzętu w wyznaczonym miejscu oraz aktualizacja jego statusu w systemie. \\
        \hline
        \textbf{Warunki wstępne} & Wyposażenie zostało zakupione i widnieje w systemie. Sprzęt fizycznie dotarł na miejsce instalacji. \\
        \hline
        \textbf{Warunki końcowe} & Wyposażenie jest gotowe do użytku, posiada przypisaną lokalizację oraz status „Dostępne”. \\
        \hline
        \textbf{Główny przepływ zdarzeń} & 
            1. Technik lokalizuje sprzęt na liście „Oczekujące na instalację”. \newline
            2. Technik wybiera odpowiedni sprzęt i przechodzi do formularza instalacji. \newline
            3. Technik wprowadza lokalizację oraz osobę odpowiedzialną. \newline
            4. Technik zmienia status sprzętu na „Dostępne”. \newline
            5. System aktualizuje stan dostępności sprzętu. \\
        \hline
        \textbf{Alternatywne przepływy zdarzeń} & 
            4a. Podczas instalacji wykryto uszkodzenie – status zmieniany jest na „Uszkodzone”. \\
        \hline
        \textbf{Specjalne wymagania} & System musi obsługiwać czytnik kodów kreskowych do identyfikacji sprzętu. \\
        \hline
        \textbf{Notatki} & brak \\
        \hline
    \end{tabular}

    \vspace{1cm}

    % --- PRZYPADEK 6: Okresowy przegląd wyposażenia ---
    \begin{tabular}{|c|c|}%{|m{10em}|m{20em}|}
        \hline
        \textbf{Nazwa} & Okresowy przegląd wyposażenia \\
        \hline
        \textbf{Numer PU} & 6 \\
        \hline
        \textbf{Aktorzy} & Serwisant \\
        \hline
        \textbf{Krótki opis} & Regularne sprawdzanie stanu technicznego sprzętu zgodnie z harmonogramem konserwacji. \\
        \hline
        \textbf{Warunki wstępne} & Nadszedł termin planowanego przeglądu (system wygenerował alert). \\
        \hline
        \textbf{Warunki końcowe} & Wpis o przeglądzie zostaje dodany do historii sprzętu. Data następnego przeglądu zostaje zaktualizowana. \\
        \hline
        \textbf{Główny przepływ zdarzeń} & 
            1. Serwisant wybiera sprzęt z listy „Do przeglądu”. \newline
            2. Serwisant przeprowadza procedurę sprawdzającą. \newline
            3. Serwisant wprowadza do systemu wynik przeglądu („Pozytywny”). \newline
            4. Serwisant ustawia datę kolejnego przeglądu. \newline
            5. System zapisuje protokół przeglądu. \\
        \hline
        \textbf{Alternatywne przepływy zdarzeń} & 
            3a. Wynik przeglądu jest „Negatywny” – system zmienia status na „Niedostępne” i generuje zlecenie naprawy. \\
        \hline
        \textbf{Specjalne wymagania} & Możliwość dołączenia skanu protokołu papierowego do rekordu w systemie. \\
        \hline
        \textbf{Notatki} & Przeglądy mogą być wymogiem gwarancyjnym producenta. \\
        \hline
    \end{tabular}

    \vspace{1cm}

    % --- PRZYPADEK 7: Rezerwacja wyposażenia ---
    \begin{tabular}{|c|c|}%{|m{10em}|m{20em}|}
        \hline
        \textbf{Nazwa} & Rezerwacja wyposażenia \\
        \hline
        \textbf{Numer PU} & 7 \\
        \hline
        \textbf{Aktorzy} & Pracownik \\
        \hline
        \textbf{Krótki opis} & Zablokowanie dostępu do sprzętu dla innych użytkowników na określony przedział czasowy. \\
        \hline
        \textbf{Warunki wstępne} & Użytkownik jest zalogowany. Wybrane wyposażenie ma status „Dostępne”. \\
        \hline
        \textbf{Warunki końcowe} & Sprzęt ma status „Zarezerwowany”. W kalendarzu sprzętu widnieje zajęty przedział czasowy. \\
        \hline
        \textbf{Główny przepływ zdarzeń} & 
            1. Pracownik wyszukuje interesujące go wyposażenie. \newline
            2. Pracownik sprawdza kalendarz dostępności. \newline
            3. Pracownik wybiera opcję „Zarezerwuj” i wskazuje termin. \newline
            4. System weryfikuje brak kolizji z innymi rezerwacjami. \newline
            5. System potwierdza rezerwację i wysyła powiadomienie. \\
        \hline
        \textbf{Alternatywne przepływy zdarzeń} & 
            4a. Konflikt terminów – system informuje, że sprzęt jest już zarezerwowany. \\
        \hline
        \textbf{Specjalne wymagania} & Czas odpowiedzi interfejsu rezerwacji poniżej 2 sekund. \\
        \hline
        \textbf{Notatki} & brak. \\
        \hline
    \end{tabular}

    \vspace{1cm}

    % --- PRZYPADEK 8: Naprawa zepsutego wyposażenia ---
    \begin{tabular}{|c|c|}%{|m{10em}|m{20em}|}
        \hline
        \textbf{Nazwa} & Naprawa zepsutego wyposażenia \\
        \hline
        \textbf{Numer PU} & 8 \\
        \hline
        \textbf{Aktorzy} & Pracownik, Technik \\
        \hline
        \textbf{Krótki opis} & Proces zgłoszenia usterki, naprawy sprzętu i przywrócenia go do stanu używalności. \\
        \hline
        \textbf{Warunki wstępne} & Wyposażenie jest uszkodzone lub nie działa poprawnie. Sprzęt znajduje się w systemie. \\
        \hline
        \textbf{Warunki końcowe} & Sprzęt ma status „Dostępne”. Historia napraw w systemie jest zaktualizowana. \\
        \hline
        \textbf{Główny przepływ zdarzeń} & 
            1. Pracownik zgłasza usterkę i zmienia status sprzętu na „Uszkodzone”. \newline
            2. System powiadamia Technika o zgłoszeniu. \newline
            3. Technik potwierdza przyjęcie zgłoszenia. \newline
            4. Technik naprawia sprzęt i wprowadza opis naprawy. \newline
            5. Technik zmienia status sprzętu na „Dostępne”. \\
        \hline
        \textbf{Alternatywne przepływy zdarzeń} & 
            4a. Naprawa niemożliwa – status zmieniany na „Do kasacji”. \newline
            4b. Wymaga zamówienia części – status zmieniany na „Oczekuje na części”. \\
        \hline
        \textbf{Specjalne wymagania} & System powinien generować raporty kosztów napraw. \\
        \hline
        \textbf{Notatki} & brak \\
        \hline
    \end{tabular}


% --- TABELA 1 ---
\noindent
\begin{tabular}{|c|c|}%{|m{10em}|m{20em}|}
        \hline
        Nazwa & Rejestracja pacjenta \\
        \hline
        Numer PU & 9\\
        \hline
        Aktorzy & Pacjent, Pracownik rejestracji\\
        \hline
        Krotki opis & Pracownik ustala wizyte z lekarzem i przekazuje informacje pacjentowi\\
        \hline
        Warunki wstepne & Pacjent zgłasza chęc umówienia wizyty\\
        \hline
        Warunki koncowe & Wizyta jest zapisana w systemie, pacjent otrzymal informacje\\
        \hline
        Glowny przepływ zdarzeń & 1. Pracownik pyta o dane pacjenta \newline
                                 2. Sprawdza wolne terminy w systemie \newline
                                 3. Rezerwuje wizyte w systemie \newline
                                 4. Przekazuje informacje pacjentowi  \\
        \hline
        Alternatywny przeplyw zdarzen &  1a. Pacjenta nie ma w systemie - pracownik dodaje nowe dane pacjenta\\
        \hline
        Specjalne wymagania & brak\\
        \hline
        Notatki& brak\\
        \hline
\end{tabular}

\vspace{1cm}

% --- TABELA 2 ---
\noindent
\begin{tabular}{|c|c|}%{|m{10em}|m{20em}|}
        \hline
        Nazwa &Wizyta lekarska \\
        \hline
        Numer PU & 10\\
        \hline
        Aktorzy & Pacjent, Lekarz\\
        \hline
        Krotki opis & pacjent konsultuje sie z lekarzem\\
        \hline
        Warunki wstepne & wizyta jest ustalona w terminarzu\\
        \hline
        Warunki koncowe & lekarz wpisal dane o wizycie do systemu\\
        \hline
        Glowny przepływ zdarzeń& 1.Pacjent przychodzi na wizyte \newline
                                 2.nastepuje konsultacja lekarska \newline
                                 3. Wizyta sie konczy  \\
        \hline
        Alternatywny przeplyw zdarzen &  2a. lekarz wystawia recepte \newline
                                       2b. lekarz wystawia skierowanie\\
        \hline
        Specjalne wymagania & brak\\
        \hline
        Notatki& brak\\
        \hline
\end{tabular}

\vspace{1cm}

% --- TABELA 3 ---
\noindent
\begin{tabular}{|c|c|}%{|m{10em}|m{20em}|}
        \hline
        Nazwa & Wystawienie recepty \\
        \hline
        Numer PU & 11\\
        \hline
        Aktorzy & Lekarz, Krajowa Baza Danych Recept\\
        \hline
        Krotki opis & Lekarz wystawia recepte i informacja jest przesylana do Bazy Danych\\
        \hline
        Warunki wstepne & Zakończona konsultacja z zaleceniem lekow\\
        \hline
        Warunki koncowe & Recepta jest zapisana w systemie i wyslana do bazy\\
        \hline
        Glowny przepływ zdarzeń & 1. Lekarz wybiera opcje wystawienia recepty \newline
                                 2. Wpisuje informacje o leku i numer lekarza \newline
                                 3. System wysyla dane do bazy  \\
        \hline
        Alternatywny przeplyw zdarzen &  3a. Brak internetu - system zapisuje recepte lokalnie\\
        \hline
        Specjalne wymagania & brak\\
        \hline
        Notatki& brak\\
        \hline
\end{tabular}

\vspace{1cm}

% --- TABELA 4 ---
\noindent
\begin{tabular}{|c|c|}%{|m{10em}|m{20em}|}
        \hline
        Nazwa & Wystawianie skierowania \\
        \hline
        Numer PU & 12\\
        \hline
        Aktorzy & Lekarz, Pacjent\\
        \hline
        Krotki opis & Lekarz wystawia skierowanie na badania lub do specjalisty\\
        \hline
        Warunki wstepne & Zdiagnozowana potrzeba dalszego leczenia\\
        \hline
        Warunki koncowe & Skierowanie jest zapisane i przekazane pacjentowi\\
        \hline
        Glowny przepływ zdarzeń & 1. Lekarz wybiera opcje skierowania \newline
                                 2. Wpisuje rodzaj badania i informacje medyczne \newline
                                 3. System generuje i zapisuje skierowanie  \\
        \hline
        Alternatywny przeplyw zdarzen &  2a. Badanie nie wymaga skierowania - lekarz informuje pacjenta\\
        \hline
        Specjalne wymagania & brak\\
        \hline
        Notatki& brak\\
        \hline
\end{tabular}

\end{document}







